% !PRINTED 506
\weberpage{506}

Im \S\,134 ist ferner nachgewiesen, daß für $D\equiv 1\ (\mathrm{mod}\ 4)$
$\sqrt{D}\,\gamma_3(\omega)$ Klasseninvariante ist. Es ist aber (\S\,54)

\weq{4}{\gamma_3(\omega) = \dfrac{8(2+\varkappa^2\varkappa'^2)(\varkappa'^2-\varkappa^2)}{\varkappa^2\varkappa'^2}.}

Genügt $\omega$ der Gleichung (2), und setzt man:

\weq{5}{\omega = \dfrac{\omega'-1}{\omega'+1},}

so ergibt sich für $\omega'$ die Gleichung:

\weq{6}{(A+B+C)\omega'^2 - 2(A-C)\omega' + (A-B+C) = 0.}

Ist $D\equiv 5\ (\mathrm{mod}\ 8)$, so sind $A$, $B$, $C$ ungerade, und ist
$D\equiv 1\ (\mathrm{mod}\ 8)$, so können wir $A$ und $C$ gerade annehmen, und
folglich ist in beiden Fällen (6) eine primitive Gleichung für $\omega'$
von der Diskriminante $4D$. Es ist daher nach \S\,126, 1. $f(\omega')^{24}$
eine Klasseninvariante für diese Diskriminante, und da nun
nach \S\,34 (18)

\[ f(\omega)f(\omega') = \sqrt{2} \]

ist, so gilt dasselbe für $f(\omega)^{24}$. Daraus folgt, daß $\varkappa^2\varkappa'^2$
Klasseninvariante für $4D$ ist, und da man nach (4) $\varkappa^2$ rational
durch $\gamma_3$ und $\varkappa^2\varkappa'^2$ ausdrücken kann, so ergibt sich:

3. \gesperrt{Ist die Diskriminante $D\equiv 1\ (\mathrm{mod}\ 4)$, so ist
$\varkappa^2$ rational ausdrückbar durch eine Klasseninvariante der Diskriminante $4D$ und durch $\sqrt{D}$.}

Ist $j(\omega)$ Klasseninvariante der Diskriminante $D$, so sind
$j(2\omega)$, $j(\frac{1}{2}\omega)$ Klasseninvarianten für $4D$, und es ist:

\weq{7}{
\begin{aligned}
j(2\omega) &= \dfrac{(f_1(2\omega)^{24}+16)^3}{f_1(2\omega)^{24}} = \dfrac{[256+f_2(\omega)^{24}]^3}{f_2(\omega)^{48}} \\
&= \dfrac{16(16\varkappa'^2+\varkappa^4)^3}{\varkappa^8\varkappa'^2}, \\
j\left(\dfrac{\omega}{2}\right) &= \dfrac{\left[f_2\!\left(\dfrac{\omega}{2}\right)^{24}+16\right]^3}{f_2\!\left(\dfrac{\omega}{2}\right)^{24}} = \dfrac{[256+f_1(\omega)^{24}]^3}{f_1(\omega)^{48}} \\
&= \dfrac{16(16\varkappa^2+\varkappa'^4)^3}{\varkappa^2\varkappa'^8}.
\end{aligned}
}

Hieraus ergibt sich, daß auch umgekehrt die Klasseninvarianten
der Diskriminante $4D$ rational durch $\varkappa^2$ ausdrückbar sind.
